\documentclass[11pt]{article}
\usepackage[margin=1in]{geometry}
\usepackage{fancyvrb}
\usepackage{booktabs}
\usepackage{hyperref}
\usepackage{parskip}
\title{bbterm --- A Beginner's Guide to Your Local Market Terminal}
\author{}
\date{}

\begin{document}
\maketitle
\tableofcontents

\section{What bbterm is}
bbterm is a market terminal you run in your own terminal window. It shows price
charts, a watchlist of symbols you follow, a statistics panel, company
fundamentals and filings, news headlines, and even congressional stock trades.
Most data comes from free sources --- Yahoo Finance (for prices) and the U.S.
Securities and Exchange Commission's EDGAR system (for fundamentals and filings)
--- and is cached locally in a small database, so the app is fast, works offline
once data is fetched, and costs almost nothing to run.

\section{Installing and starting the app}
The easiest way to install is with \texttt{pipx} (which requires Python 3.12 or
newer):
\begin{Verbatim}[frame=single]
pipx install bbterm-tui      # the command is still called "bbterm"
bbterm                       # launch the terminal
\end{Verbatim}
If you are working from a clone of the source instead, run \texttt{bbterm} from
the project's virtual environment (\texttt{.venv/bin/bbterm}).

bbterm runs on free data out of the box --- no keys required. Two optional keys
unlock extra features, set as environment variables (or in a \texttt{.env} file
in the folder you run bbterm from):
\begin{center}
\begin{tabular}{ll}
\toprule
Variable & Unlocks \\
\midrule
\texttt{DATABENTO\_API\_KEY} & Higher-quality price history from Databento \\
\texttt{LAMBDA\_API\_KEY} & Congressional trades (the \texttt{POL} view) \\
\bottomrule
\end{tabular}
\end{center}

\textbf{Best experience:} run bbterm in an image-capable terminal such as
\textbf{iTerm2}, where the price chart is drawn as a crisp picture. In a plain
terminal, bbterm automatically falls back to a text chart.

\section{The screen}
\begin{Verbatim}[frame=single]
+----------------------------------------------+
| : AAPL_                          (command)   |  <- command bar (top)
+----------+-----------------------------------+
| WATCHLIST|  AAPL  291.52  -6.48 (-2.18%)     |
| SPY  ... |  [ chart / stats / FA / FIL /     |  <- main panel
| AAPL ... |    news / congress ]              |
+----------+-----------------------------------+
| SPY .. | QQQ .. | AAPL ..  (ticker strip)    |
| [GP]chart [DES]stats [FA][FIL][N][POL] ?=help|  <- footer
+----------------------------------------------+
\end{Verbatim}
\textbf{Watchlist} (left): the symbols you follow, each with price and percent
change. \textbf{Main panel} (right): whichever view you have chosen. \textbf{Ticker
strip} (bottom): a compact running line of your watchlist.

\section{The command bar and the colon key}
The command bar is one line at the top where you type instructions instead of
hunting through menus. Because single keys are also shortcuts (for example
\texttt{q} quits), the bar is normally ``asleep''. Press \texttt{:} (the colon
key) to wake it up and move your typing into it. Type your command, press
\texttt{Enter} to run it, or press \texttt{Esc} to leave the bar without running
anything. This colon-to-type idea is borrowed from the \texttt{vim} editor.

\section{Commands}
\begin{center}
\begin{tabular}{ll}
\toprule
Type this & What happens \\
\midrule
\texttt{AAPL} & Load Apple into the main panel \\
\texttt{ADD TSLA} & Add Tesla to your watchlist (saved) \\
\texttt{DEL SPY} & Remove SPY from your watchlist (saved) \\
\texttt{GP} & Show the price chart \\
\texttt{DES} & Show the statistics panel \\
\texttt{FA} & Show company fundamentals (from SEC EDGAR) \\
\texttt{FIL} & Show recent SEC filings \\
\texttt{N} & Show recent news headlines \\
\texttt{POL} & Show congressional trades for this stock \\
\texttt{?} or \texttt{HELP} & Show a reminder of the commands \\
\bottomrule
\end{tabular}
\end{center}

\textbf{Ticker tip:} for ``Class B''-style symbols, use a hyphen, not a dot ---
for example Berkshire Hathaway Class B is \texttt{BRK-B} (typing \texttt{BRKB} or
\texttt{BRK.B} returns no data on the free source).

\section{Keyboard shortcuts}
These work when the command bar is \emph{not} focused.
\begin{center}
\begin{tabular}{ll}
\toprule
Key & Action \\
\midrule
\texttt{:} & Focus the command bar \\
\texttt{1}--\texttt{6} & Time range: 1D, 5D, 1M, 6M, 1Y, 5Y \\
\texttt{r} & Refresh data \\
\texttt{q} & Quit \\
\bottomrule
\end{tabular}
\end{center}

\section{Reading the chart}
The price chart (\texttt{GP}) is a candlestick chart. Each candle is one time
period: the thick body spans the open and close prices --- green if the price
rose, red if it fell --- and the thin lines (wicks) above and below show the
highest and lowest prices reached. The smaller panel beneath shows trading
volume. Use the number keys \texttt{1}--\texttt{6} to change the time range. In an
image-capable terminal (such as iTerm2) the chart is a sharp picture; elsewhere
it is drawn with text characters.

\section{Fundamentals and filings (FA / FIL)}
Type \texttt{FA} to see headline financials --- revenue, net income, earnings
per share, assets, and more --- with each figure's fiscal year and its change
versus the prior year. Type \texttt{FIL} to see the company's most recent filings
with the U.S. Securities and Exchange Commission (annual 10-K, quarterly 10-Q,
and event-driven 8-K reports). The filings view is a list: use the up and down
arrow keys to highlight a filing and press \texttt{Enter} to open it in your web
browser. This data comes from SEC EDGAR, is free and public, and is cached on
your machine for a day.

\section{News (N)}
Type \texttt{N} to see recent news headlines for the loaded symbol, each with its
source, how long ago it was published, and a link to the article. Headlines come
from a free news feed and are cached briefly so the view stays fast. (To open an
article, copy its link.)

\section{Congressional trades (POL)}
Type \texttt{POL} to see recent stock trades in the loaded symbol made by a
curated list of well-known members of Congress (such as Nancy Pelosi). For each
person you get a short summary --- how many buys, how many sells, and an estimated
net dollar amount --- followed by the individual disclosed trades.

A few honest limits to keep in mind:
\begin{itemize}
  \item This uses official congressional disclosures, which report \emph{dollar
  ranges} (for example ``\$15,001 -- \$50,000''), never exact share counts.
  \item It shows \emph{trades} (buys and sells), not a person's current holdings.
  \item It covers \emph{members of Congress only} --- not the President, the Vice
  President, their families, or private individuals, who do not file these
  reports.
\end{itemize}
This view needs a free \texttt{LAMBDA\_API\_KEY}; without one, bbterm shows a short
note telling you how to enable it.

\section{Reading the statistics panel (DES)}
\begin{center}
\begin{tabular}{ll}
\toprule
Field & Meaning \\
\midrule
Last & Most recent price \\
Change & Price and percent move versus the previous close \\
52w High / Low & Highest and lowest price over the past year \\
1M Return & Percent change over about one month \\
YTD Return & Percent change since January 1st \\
Avg Vol & Average daily trading volume \\
Day Range & Today's low and high \\
\bottomrule
\end{tabular}
\end{center}
All statistics are computed from the price history already on your machine.

\end{document}
